\documentclass[11pt]{amsart}

\usepackage[T1]{fontenc}
\usepackage{lmodern}
\usepackage{microtype}
\usepackage{mathtools,amssymb,amsthm}
\usepackage[hidelinks]{hyperref}

\hypersetup{
  pdftitle={Eliminating the 4r-7 Alternative for Diameter-Two Mixed Almost Moore Graphs},
  pdfsubject={Mixed almost Moore graphs of diameter two},
  pdfkeywords={mixed graph, almost Moore graph, degree-diameter problem, adjacency matrix, trace}
}

\newtheorem{theorem}{Theorem}
\newtheorem{corollary}[theorem]{Corollary}

\DeclareMathOperator{\tr}{tr}

\title[Eliminating the $4r-7$ alternative]
{Eliminating the \texorpdfstring{$4r-7$}{4r-7} Alternative for Diameter-Two Mixed Almost Moore Graphs}
\date{}

\subjclass[2020]{05C35, 05C50}
\keywords{mixed graph, almost Moore graph, degree--diameter problem, adjacency matrix, trace}

\begin{document}

\begin{abstract}
L\'opez and Miret derived two arithmetic alternatives for the
parameters of a totally regular mixed almost Moore graph of diameter
two. We prove that their alternative $c^{2}=4r-7$ cannot occur when
the undirected degree satisfies $r>2$ and the directed degree
satisfies $z\geq 1$. Consequently, in the standard simple mixed-graph
model, every such graph with $r>2$ and $z\geq 1$ must satisfy
\[
  c^{2}=4r+1,
  \qquad
  c\mid(4z+1)(4z-7)
\]
for some odd integer $c$. The square condition is equivalent to
$r=s(s+1)$ for an integer $s\geq 2$. In particular, no such graph has
undirected degree $r=4$, settling the first open row in the table of
L\'opez and Miret.
\end{abstract}

\maketitle

\section{Introduction}

A mixed almost Moore graph of diameter two and degrees $(r,z)$ has
order
\[
  n=(r+z)^2+z.
\]
For totally regular graphs with even undirected degree $r>2$ and
directed degree $z\geq 1$, L\'opez and Miret proved that one of the
following alternatives is necessary \cite[Theorem~2]{LM}:
\begin{align}
  c_1^2&=4r+1,
  & c_1&\mid(4z+1)(4z-7), \label{eq:first-branch}\\
  c_2^2&=4r-7,
  & c_2&\mid(16z^2+40z-23). \label{eq:second-branch}
\end{align}
Here $c_1,c_2$ are odd integers. The family $r=4$, which lies in
\eqref{eq:second-branch}, was listed as unknown in
\cite[Table~1]{LM} and again in \cite[Table~1]{DFL}.

In the proof of \cite[Proposition~4]{LM}, the cubic-trace computation
is displayed for \eqref{eq:first-branch}, while the other branch is
said to follow by the same method. We carry out the branch-dependent
calculation. The relevant difference of cubes is $c(r-1)$ in
\eqref{eq:second-branch}, rather than $c(r+1)$ as in
\eqref{eq:first-branch}; the resulting trace is negative, which is
impossible.

\section{The trace obstruction}

\begin{theorem}\label{thm:main}
There is no totally regular mixed almost Moore graph of diameter two
with undirected degree $r>2$, directed degree $z\geq 1$, and $4r-7$
a square.
\end{theorem}

\begin{proof}
Let $d=r+z$ and
\[
  n=d^2+d-r.
\]
Let $A$ be the adjacency matrix of the graph, let $P$ be the
permutation matrix of its repeat automorphism, and let $J$ be the
all-ones matrix. By \cite[Eq.~(1)]{LM},
\begin{equation}\label{eq:matrix}
  I+A+A^2=J+rI+P.
\end{equation}
Moreover, total regularity and invariance under the repeat
automorphism give
\begin{equation}\label{eq:commutation}
  AJ=JA=dJ,
  \qquad
  AP=PA.
\end{equation}
Put
\[
  f=\tr P,
  \qquad
  t=\tr(AP).
\]
Since $A$ and $P$ are nonnegative integer matrices, both are
nonnegative integers. Take traces in \eqref{eq:matrix}, and then
multiply \eqref{eq:matrix} by $A$ and take traces again. Using
$\tr A=0$, $\tr J=n$, and $\tr(AJ)=d\tr J=dn$, the last by
\eqref{eq:commutation}, this gives
\begin{equation}\label{eq:combinatorial-traces}
  \tr(A^2)=rn+f,
  \qquad
  \tr(A^3)=zn-f+t.
\end{equation}

Assume that $c^2=4r-7$ with $c>0$. Since $4r-7$ is odd, so is $c$,
whence $c^2\equiv1\pmod 8$ and
$r$ is even; together with $r>2$ and $z\geq 1$ this places the
graph under the hypotheses of \cite[Theorem~2]{LM}. Moreover, the
integer $4r+1$ is not a square: otherwise, if $h^2=4r+1$ with $h>0$,
then $h$ is odd, again because $4r+1$ is odd, and
\[
  (h-c)(h+c)=8,
\]
and since $h$ and $c$ are odd both factors are even, while
$h^2-c^2=8>0$ makes both positive, so $\{h-c,h+c\}=\{2,4\}$. This
forces $h=3$, $c=1$, and $r=2$, contradicting $r>2$. Thus the third
case in the proof of \cite[Theorem~2]{LM} applies.

Let $q_j$ be the number of cycles of length $j$ in $P$, and define
\[
  s_i=\sum_{i\mid j}q_j.
\]
Euler's totient function $\varphi$ and the M\"obius function $\mu$
satisfy
\begin{equation}\label{eq:mobius-identities}
  \sum_{i=1}^{n}s_i\varphi(i)=n,
  \qquad
  \sum_{i=1}^{n}s_i\mu(i)=q_1=f.
\end{equation}
Indeed, these follow by summing the identities
\[
  \sum_{i\mid j}\varphi(i)=j,
  \qquad
  \sum_{i\mid j}\mu(i)=\mathbf{1}_{j=1}
\]
over the cycles of $P$.

Set
\[
  \alpha=\frac{-1+c}{2},
  \qquad
  \beta=\frac{-1-c}{2}.
\]
The characteristic-polynomial factorization in the third case of
\cite[proof of Theorem~2]{LM} is
\begin{equation}\label{eq:charpoly}
\begin{split}
  \chi_A(x)={}&(x-d)(x^2+x-r)^{(s_1-1)/2}
  (x-\alpha)^a(x-\beta)^b\\
  &\times
  \prod_{i=3}^{n}
  \Phi_i\!\left(x^2+x-(r-1)\right)^{s_i/2},
\end{split}
\end{equation}
where $a,b\in\mathbb{Z}_{\geq 0}$, $a+b=s_2$, and $\Phi_i$ is the
$i$th cyclotomic polynomial. The branch enters at the two lowest
indices. Since $\Phi_1(y)=y-1$ and $\Phi_2(y)=y+1$, the polynomials
$\Phi_i\!\left(x^2+x-(r-1)\right)$ for $i=1$ and $i=2$ are
$x^2+x-r$ and $x^2+x-(r-2)$, of discriminants $4r+1$ and $4r-7$.
Under \eqref{eq:second-branch} it is the second that splits over
$\mathbb{Z}$, its roots being $\alpha$ and $\beta$, while the first
stays irreducible; under \eqref{eq:first-branch} the roles are
exchanged. This is the source of the $c(r-1)$ in
\eqref{eq:branch-cubes} below, in place of the $c(r+1)$ of the first
branch.

Write $\Delta=a-b$. Taking the sum of the roots in
\eqref{eq:charpoly} and using
\eqref{eq:mobius-identities}, we obtain
\[
  0=\tr A
  =d-\frac{s_1-1}{2}+\frac{-s_2+c\Delta}{2}
   -\frac12\sum_{i=3}^{n}s_i\varphi(i)
  =d+\frac{1+c\Delta-n}{2}.
\]
Hence
\begin{equation}\label{eq:delta}
  c\Delta=n-2d-1.
\end{equation}

If $u,v$ are the roots of $x^2+x-r$, then
\[
  u^3+v^3=-1-3r.
\]
Moreover,
\begin{equation}\label{eq:branch-cubes}
  \alpha^3+\beta^3=5-3r,
  \qquad
  \alpha^3-\beta^3=c(r-1).
\end{equation}
For a primitive $i$th root of unity $\zeta$, let
$\gamma_\zeta,\delta_\zeta$ be the roots of
\[
  x^2+x-(r-1)=\zeta.
\]
Then
\[
  \gamma_\zeta^3+\delta_\zeta^3=2-3r-3\zeta.
\]
Summing the cubes of the roots in \eqref{eq:charpoly}, evaluating
\[
  2\bigl(a\alpha^3+b\beta^3\bigr)
  =s_2(\alpha^3+\beta^3)+\Delta(\alpha^3-\beta^3)
\]
by \eqref{eq:branch-cubes}, and using \eqref{eq:mobius-identities},
yields
\begin{align}
  2\bigl(\tr(A^3)-d^3\bigr)
  ={}&(-1-3r)(s_1-1)+(5-3r)s_2+(r-1)c\Delta
  \notag\\
  &+\sum_{i=3}^{n}s_i
  \bigl((2-3r)\varphi(i)-3\mu(i)\bigr)
  \notag\\
  ={}&(r-1)c\Delta+(2-3r)n+3r+1-3f.
  \label{eq:spectral-cubic}
\end{align}

Comparing \eqref{eq:combinatorial-traces} and
\eqref{eq:spectral-cubic}, and then using \eqref{eq:delta},
$z=d-r$, and $n=d^2+d-r$, gives
\begin{align}
  2t
  &={}2d^3-2zn+(r-1)(n-2d-1)
       +(2-3r)n+3r+1-f
  \notag\\
  &=-d^2+3d+r+2-f.
  \label{eq:t}
\end{align}

Recall that $r$ is even; with $r>2$ this gives $r\geq 4$, and
$z\geq 1$ gives $d=r+z\geq r+1$. As $f\geq 0$ and $-d^2+3d$ is
decreasing for $d\geq 2$, \eqref{eq:t} gives
\[
  2t
  \leq -(r+1)^2+3(r+1)+r+2
  =-r^2+2r+4
  <0,
\]
the final inequality because $r\geq 4$. This contradicts $t\geq 0$.
\end{proof}

\begin{corollary}\label{cor:parameters}
In the standard simple mixed-graph model, every mixed almost Moore
graph of diameter two with $r>2$ and $z\geq 1$ satisfies
\[
  c^2=4r+1,
  \qquad
  c\mid(4z+1)(4z-7)
\]
for some odd integer $c$. Equivalently, the square condition can be
written as
\[
  r=s(s+1)
\]
for an integer $s\geq 2$. In particular, no such graph has $r=4$,
for any $z\geq 1$.
\end{corollary}

\begin{proof}
Mixed almost Moore graphs of diameter two are totally regular
\cite[Theorem~2.8]{TE}, and their undirected degree is even
\cite[Proposition~1]{LM}. Thus the hypotheses of
\cite[Theorem~2]{LM} are met, $r>2$ and $z\geq 1$ by assumption and
the parity of $r$ by the foregoing, and that theorem applies.
Theorem~\ref{thm:main} eliminates its second alternative, leaving the
first. Writing the positive odd integer $c$ as $2s+1$ gives
\[
  r=\frac{c^2-1}{4}=s(s+1).
\]
For $r=4$, the remaining alternative would require $c^2=17$, which
is impossible.
\end{proof}

\section{Status and scope}

Corollary~\ref{cor:parameters} gives a partial negative answer to
\cite[Question~3]{LM} and to the broader diameter-two existence
problem \cite[Problem~5.2]{Shallcross}. It settles the $r=4$ family
that was recorded as unknown in \cite[Table~1]{LM} and
\cite[Table~1]{DFL}.

What is settled is one of the two arithmetic alternatives of
\cite[Theorem~2]{LM}: Theorem~\ref{thm:main} refutes the branch
$c^2=4r-7$ for every $r>2$ and every $z\geq 1$, uniformly in both
degrees, so the $r=4$ row is a special case of a general statement
rather than a finite check. What is not settled is the branch
$c^2=4r+1$; for the degrees $r=s(s+1)$ with $s\geq 2$ the existence
question remains open. The hypotheses $r>2$ and $z\geq 1$ are those
of \cite[Theorem~2]{LM}, whose remaining hypothesis that $r$ be even
is automatic here, and the hypothesis $z\geq 1$ is used rather than
decorative: substituting $z=0$, that is $d=r$, in \eqref{eq:t} gives
$2t=-r^2+4r+2-f$, which at $r=4$ is $2-f$ and so permits $t\geq 0$.
The argument therefore says nothing about $z=0$, nor about $r\leq 2$.

The proof uses no computation: its decisive step is the evaluation
of \eqref{eq:t} at $d=r+1$, which can be checked by hand. It is not
self-contained. The matrix identity \eqref{eq:matrix} is quoted from
\cite[Eq.~(1)]{LM} and the characteristic-polynomial factorization
\eqref{eq:charpoly} from \cite[proof of Theorem~2]{LM}, whose case
analysis we do not reproduce; Corollary~\ref{cor:parameters} further
imports total regularity from \cite[Theorem~2.8]{TE} and the parity
of $r$ from \cite[Proposition~1]{LM}. What is carried out here is the
branch-dependent cubic-trace calculation and the consequence drawn
from it. We are not aware of prior work eliminating the
$4r-7$ branch; the literature we have located that cites \cite{LM},
including \cite{DFL}, restates the $r=4$ case as open.

\begin{thebibliography}{9}

\bibitem{DFL}
C.~Dalf\'o, M.~A. Fiol, and N.~L\'opez,
\emph{Almost Moore and the largest mixed graphs of diameters two and
three},
Linear Algebra Appl. \textbf{693} (2024), 374--385.
\href{https://doi.org/10.1016/j.laa.2024.01.007}
{doi:10.1016/j.laa.2024.01.007}.

\bibitem{LM}
N.~L\'opez and J.~M. Miret,
\emph{On mixed almost Moore graphs of diameter two},
Electron. J. Combin. \textbf{23} (2016), no.~2,
Paper P2.3, 14~pp.
\href{https://doi.org/10.37236/5647}
{doi:10.37236/5647}.

\bibitem{Shallcross}
E.~Shallcross,
\emph{On the total regularity of almost mixed Moore graphs},
\href{https://arxiv.org/abs/2608.12425}
{arXiv:2608.12425 [math.CO]}, 2026.

\bibitem{TE}
J.~Tuite and G.~Erskine,
\emph{On total regularity of mixed graphs with order close to the
Moore bound},
Graphs Combin. \textbf{35} (2019), 1253--1272.
\href{https://doi.org/10.1007/s00373-019-02114-2}
{doi:10.1007/s00373-019-02114-2}.

\end{thebibliography}

\end{document}